\documentclass[12pt]{article}
\usepackage{amsmath,amssymb,amsthm,hyperref}

\newtheorem{theorem}{Theorem}
\newtheorem{proposition}[theorem]{Proposition}
\newtheorem{lemma}[theorem]{Lemma}
\newtheorem{corollary}[theorem]{Corollary}
\theoremstyle{definition}
\newtheorem{definition}[theorem]{Definition}
\newtheorem{hyp}[theorem]{Hypothesis}
\newtheorem{remark}[theorem]{Remark}

\title{On Runs of Integers with Large Prime Factors:\\
A Lower Bound for $k(n)$ and the Status of the Upper Bound}
\author{}
\date{April 2026}

\begin{document}
\maketitle

\begin{abstract}
Let $k(n)$ be the maximal $k$ such that there exists $m \leq n$ with each of
$m+1, \ldots, m+k$ divisible by some prime $> k$.
We prove unconditionally that
\[
  k(n) \;\geq\; \exp\!\bigl((\log n)^{1/2 - \varepsilon}\bigr)
\]
for every fixed $\varepsilon > 0$ and all sufficiently large $n$,
using Rankin's method and a pigeonhole argument.

We also discuss the upper bound $k(n) \leq \exp((\log n)^{1/2+\varepsilon})$,
which would give $\log k(n) = (\log n)^{1/2+o(1)}$.
The upper bound argument requires the existence of a $k_1$-smooth integer
in every interval $[m+1, m+k_1]$ with $m \leq n$ and $k_1 = \exp((\log n)^{1/2+\varepsilon})$.
We reduce this to a short-interval smooth number estimate
(Hypothesis~\ref{hyp:SI} below) and prove the upper bound
\emph{conditional on that hypothesis}.
We discuss what the available literature (Hildebrand--Tenenbaum \cite{HT},
Tenenbaum \cite{Ten}) actually establishes and where the gap lies.
\end{abstract}

\section{Setup and Definitions}

\begin{definition}
A positive integer $m$ is \emph{$y$-smooth} if every prime factor of $m$
is $\leq y$.  Write $P^+(m)$ for the largest prime factor of $m$
(with $P^+(1) = 1$ by convention).  Set
\[
  \Psi(x, y) := \#\{m \leq x : P^+(m) \leq y\}.
\]
\end{definition}

\begin{definition}
A \emph{$k$-run below $n$} is an interval $\{m+1, \ldots, m+k\}$
with $m \leq n$ and $P^+(m+j) > k$ for all $1 \leq j \leq k$.
Define $k(n)$ to be the maximum $k$ for which a $k$-run below $n$ exists.
\end{definition}

\section{The Rankin Bound}

\begin{lemma}[Rankin bound {\cite[Theorem~III.5.1]{Ten}}]\label{lem:rankin}
There is an effectively computable absolute constant $A \geq 3$ such that
for all $x \geq y \geq 2$ with $v := \log x / \log y \geq A$,
\[
  \Psi(x, y) \;\leq\; x \cdot e^{-\frac{1}{2}v\log v}.
\]
\end{lemma}
\begin{proof}
Cited from \cite[Theorem~III.5.1]{Ten}.
Rankin's method with the choice $\sigma = 1 - (\log v)/\log y$ yields
$\Psi(x,y) \leq x\,e^{-v\log v + C_0 v/\log v}$ for an absolute constant $C_0$.
The correction term $C_0 v/\log v \leq \frac{1}{2}v\log v$ once $\log^2 v \geq 2C_0$,
which holds for all $v \geq A$ with $A = e^{\sqrt{2C_0}}$.
The full proof with all constants is in \cite[Theorem~III.5.1]{Ten}.
\end{proof}

\section{The Lower Bound}

\begin{theorem}[Lower bound]\label{thm:lower}
For every fixed $\varepsilon > 0$ and all sufficiently large $n$,
$k(n) \geq \exp\!\bigl((\log n)^{1/2-\varepsilon}\bigr)$.
\end{theorem}

\begin{proof}
Set $k_0 := \lfloor\exp((\log n)^{1/2-\varepsilon})\rfloor$.

\medskip
\noindent\textit{Step 1: $k_0$-smooth numbers are sparse in $[1,n]$.}
Set $v_0 := \log n / \log k_0$.
Since $\log k_0 = (\log n)^{1/2-\varepsilon}(1+o(1))$,
we get $v_0 = (\log n)^{1/2+\varepsilon}(1+o(1)) \to \infty$.
In particular $v_0 \geq A$ (the constant from Lemma~\ref{lem:rankin})
for all large $n$, so
\[
  \Psi(n, k_0) \;\leq\; n \cdot e^{-\frac{1}{2}v_0 \log v_0}.
\]
We estimate the exponent:
\[
  \tfrac{1}{2}v_0 \log v_0
  = \tfrac{1}{2}(\log n)^{1/2+\varepsilon}(1+o(1)) \cdot (1/2+\varepsilon)\log\log n\,(1+o(1))
  \;\geq\; (\log n)^{1/2+\varepsilon/2}
\]
for all large $n$.  Since $\log k_0 = (\log n)^{1/2-\varepsilon}$, this gives
\[
  e^{-\frac{1}{2}v_0\log v_0} \leq e^{-(\log n)^{1/2+\varepsilon/2}} \ll n^{-1} \cdot k_0^{-1},
\]
so in particular $\Psi(n, k_0) \leq n/(2k_0)$ for all sufficiently large $n$.

\medskip
\noindent\textit{Step 2: Pigeonhole gives an empty block.}
Let $B := \lfloor n/k_0 \rfloor$.
Partition $\{1, \ldots, B k_0\}$ into $B$ blocks $I_i := [(i-1)k_0+1, ik_0]$.
Every $k_0$-smooth integer in $\{1,\ldots,n\}$ lies in one of these blocks, and
\[
  \Psi(n, k_0) \;\leq\; \frac{n}{2k_0} \;\leq\; \frac{B+1}{2} \;<\; B
\]
for all $n \geq 2k_0$ (i.e.\ all large $n$).  By pigeonhole,
some block $I_i$ contains no $k_0$-smooth integer.

\medskip
\noindent\textit{Step 3: The empty block is a $k_0$-run.}
For that $i$, every integer in $I_i$ has largest prime factor $> k_0$.
Setting $m = (i-1)k_0 \leq n$, the interval $\{m+1,\ldots,m+k_0\}$
is a $k_0$-run below $n$, so
$k(n) \geq k_0 = \exp((\log n)^{1/2-\varepsilon}(1+o(1)))$.
\end{proof}

\section{The Upper Bound: Reduction and Gap}

\subsection*{What needs to be proved}

Set $k_1 := \lfloor\exp((\log n)^{1/2+\varepsilon})\rfloor$.
The upper bound $k(n) \leq k_1$ is equivalent to: every interval $[m+1, m+k_1]$
with $m \leq n$ contains at least one $k_1$-smooth integer.

For $m < k_1$ this is immediate: $k_1 \in [m+1, m+k_1]$ and $k_1$ is $k_1$-smooth.

For $m \in [k_1, n]$, we need the following estimate.

\begin{hyp}[Short-interval smooth numbers]\label{hyp:SI}
Let $\rho$ denote the Dickman function and let $\varepsilon > 0$ be fixed.
Set $k_1 := \lfloor\exp((\log n)^{1/2+\varepsilon})\rfloor$.
There exists $n_* = n_*(\varepsilon)$ such that for all $n \geq n_*$
and all integers $m \in [k_1, n]$,
\begin{equation}\label{eq:SI}
  \Psi(m + k_1,\, k_1) \;-\; \Psi(m,\, k_1) \;\geq\; 1.
\end{equation}
(Since $m$ is a positive integer, the difference counts exactly the $k_1$-smooth
integers in the interval $(m, m+k_1] = \{m+1,\ldots,m+k_1\}$.)
\end{hyp}

\begin{theorem}[Upper bound, conditional]\label{thm:upper}
Assume Hypothesis~\ref{hyp:SI}.
For every fixed $\varepsilon > 0$ and all sufficiently large $n$,
$k(n) \leq \exp((\log n)^{1/2+\varepsilon})$.
\end{theorem}

\begin{proof}
With Hypothesis~\ref{hyp:SI} every interval $[m+1, m+k_1]$ with $m \in [k_1,n]$
contains a $k_1$-smooth integer, and the case $m < k_1$ was handled above.
Hence no $k_1$-run below $n$ exists, so $k(n) < k_1$.
\end{proof}

\subsection*{Discussion: the subtraction approach, its failure, and two routes to a proof}

The main result of Hildebrand--Tenenbaum \cite{HT} is a uniform asymptotic
for the \emph{global count}.  Writing
$\mathcal{R}_\varepsilon = \{(x,y) : x \geq y \geq 2,\, \log y \geq (\log\log x)^{5/3+\varepsilon}\}$
and $\eta(x,y) := \Psi(x,y)/(x\rho(u)) - 1$ with $u = \log x/\log y$,
\cite[Th\'eor\`eme~1]{HT} gives
\begin{equation}\label{eq:HTglobal}
  \eta(x,y) \to 0 \quad\text{uniformly for }(x,y)\in\mathcal{R}_\varepsilon
  \text{ as }y\to\infty.
\end{equation}
Explicitly: for every $\delta > 0$ there exists $Y_0(\varepsilon,\delta)$
(independent of $x$) such that $|\eta(x,y)| \leq \delta$ for all
$(x,y)\in\mathcal{R}_\varepsilon$ with $y \geq Y_0$.

\subsubsection*{A natural case split at $m = k_1^2$}

The range $m \in [k_1, n]$ splits naturally at $m = k_1^2$,
equivalently at $u_1 := \log m/\log k_1 = 2$.

\noindent\textbf{Sub-range~1a}: $m \in [k_1, k_1^2)$, i.e.\ $u_1 \in [1,2)$.
Write $m = qk_1 + r$ with $0 \leq r < k_1$ (Euclidean division).
The integer $(q+1)k_1 = m - r + k_1$ satisfies $m < (q+1)k_1 \leq m + k_1$,
so $(q+1)k_1 \in [m+1, m+k_1]$.
Since $m < k_1^2$ we have $q \leq k_1 - 1$, giving $q+1 \leq k_1$.
Every prime factor of $q+1$ is at most $q+1 \leq k_1$, so $(q+1)k_1$ is $k_1$-smooth.
Thus every interval $[m+1, m+k_1]$ with $m \in [k_1, k_1^2)$
contains a $k_1$-smooth integer by an elementary divisibility argument,
with no analytic input.

\noindent\textbf{Sub-range~1b}: $m \in [k_1^2, n]$, i.e.\ $u_1 \geq 2$.
Here we must use smooth-number asymptotics.
We analyze the subtraction approach for this range.

\subsubsection*{The subtraction approach for sub-range~1b and why it fails}

Subtracting \eqref{eq:HTglobal} at $m + k_1$ and at $m$:
\begin{equation}\label{eq:diff}
  \Psi(m+k_1,k_1) - \Psi(m,k_1)
  \;=\; \underbrace{(m+k_1)\rho(u') - m\rho(u_1)}_{\text{(I)}}
  \;+\; \underbrace{(m+k_1)\rho(u')\eta_1 - m\rho(u_1)\eta_2}_{\text{(II)}},
\end{equation}
where $u' = \log(m+k_1)/\log k_1$ and $\eta_j$ is $\eta(\cdot)$ at the
respective argument.

\emph{Term~(I) for $u_1 \geq 2$.}
Set $g(t) = t\rho(\log t/\log k_1)$.  By the mean-value theorem,
$g(m+k_1) - g(m) = k_1 g'(\xi)$ for some $\xi \in [m, m+k_1]$.
Computing the derivative:
\[
  g'(t) = \rho(u_t)\Bigl(1 - \frac{\rho(u_t-1)}{u_t\,\rho(u_t)\,\log k_1}\Bigr),
  \qquad u_t := \frac{\log t}{\log k_1}.
\]
For $u_t \geq 2$, the Dickman delay-differential equation $\rho'(u) = -\rho(u-1)/u$
and the standard asymptotics \cite[Corollary~III.5.4]{Ten} give
$\rho(u-1)/\rho(u) = O(u\log u)$ for $u \geq 2$.
The relative correction is therefore
\[
  \frac{\rho(u_t-1)}{u_t\,\rho(u_t)\,\log k_1}
  = O\!\left(\frac{\log u_1}{\log k_1}\right)
  \leq \frac{(1/2-\varepsilon)\log\log n}{(\log n)^{1/2+\varepsilon}} \to 0.
\]
Hence, for sub-range~1b:
\begin{equation}\label{eq:termI}
  \text{(I)} = k_1\rho(u_1)\!\left(1 + O\!\left(\frac{\log u_1}{\log k_1}\right)\right)
  = k_1\rho(u_1)(1+o(1)).
\end{equation}

\begin{remark}
The estimate \eqref{eq:termI} relies on $\rho(u_t-1)/\rho(u_t) = O(u_t\log u_t)$,
which holds for $u_t \geq 2$ but \emph{fails} for $u_t \in (1,2)$.
For $u \in (1,2)$, $\rho(u) = 1-\log u$ and $\rho(u-1) = 1$,
so $\rho(u-1)/\rho(u) = 1/(1-\log u) \to \infty$ as $u\to 1^+$.
This is why sub-range~1a requires the separate elementary argument above.
\end{remark}

\emph{Term~(II) is the fundamental obstacle.}
With $|\eta_1|, |\eta_2| \leq \delta$:
\[
  |\text{(II)}| \;\leq\;
  \bigl[(m+k_1)\rho(u') + m\rho(u_1)\bigr]\delta
  \;\lesssim\; 2m\rho(u_1)\delta.
\]
The ratio to the main term is
\[
  \frac{|\text{(II)}|}{k_1\rho(u_1)} \;\lesssim\; \frac{2m\delta}{k_1}.
\]
For $m \in [k_1^2, n]$ this ranges from $2k_1\delta$ to $2(n/k_1)\delta$.
Since $n/k_1 = \exp(\log n - (\log n)^{1/2+\varepsilon}) \to \infty$,
even at the lower end of sub-range~1b the ratio is $2k_1\delta \to \infty$
for any fixed $\delta > 0$.
No choice of $\delta$ depending only on $\varepsilon$ can control this:
\emph{the error swamps the main term throughout sub-range~1b.}

\textbf{Conclusion}: The subtraction of two applications of \eqref{eq:HTglobal}
does not prove Hypothesis~\ref{hyp:SI} for sub-range~1b.

\subsubsection*{Route~1: A Lipschitz bound on $\eta$}

If one could establish, for $(x,y)\in\mathcal{R}_\varepsilon$ and $y\geq Y_1(\varepsilon)$,
\begin{equation}\label{eq:LipHT}
  \bigl|\eta(x+y,y) - \eta(x,y)\bigr| \;\leq\; \frac{Cy}{x\log y}
\end{equation}
for an absolute constant $C$, then in \eqref{eq:diff}:
\[
  |\text{(II)}|
  \leq (m+k_1)\rho(u')|\eta_1 - \eta_2| + k_1\rho(u_1)|\eta_2|
  \leq \frac{2Ck_1\rho(u_1)}{\log k_1} + k_1\rho(u_1)\delta.
\]
Choosing $\delta = 1/4$ and $k_1 \geq e^{8C}$ (ensured by $n$ large),
Term~(II) is at most $\tfrac12 k_1\rho(u_1)$.
Since $k_1\rho(u_1) \geq 2$ (the main-term lower bound, established by the
same calculation as in the lower bound proof but with $k_1$ in place of $k_0$),
the difference is $\geq 1$.

Bound~\eqref{eq:LipHT} says the relative error $\eta$ is Lipschitz in $x$
with constant $O((x\log y)^{-1})$.
This is the equicontinuity expected from the saddle-point analysis
(the saddle point $\sigma_0 = 1 - (\log u)/\log y$ shifts by $O(1/(x\log y))$
as $x$ increases by $y$), but establishing it rigorously requires a
second-order expansion of $\Psi(x,y)$ in $1/\log y$ that is \emph{uniform in $x$}.
This is not stated in \cite{HT,Ten}.

\subsubsection*{Route~2: Direct Perron analysis}

Rather than subtracting two global estimates, one can work with
the combined Perron integral for the difference directly:
\begin{equation}\label{eq:PerronSI}
  \Psi(m+k_1,k_1) - \Psi(m,k_1)
  = \frac{1}{2\pi i}\int_{\sigma_0 - iT}^{\sigma_0+iT}
    F_{k_1}(s)\,\frac{(m+k_1)^s - m^s}{s}\,ds + \mathcal{E}(m,T),
\end{equation}
where $F_y(s) = \prod_{p\leq y}(1-p^{-s})^{-1}$,
$\sigma_0 = 1 - (\log u_1)/\log k_1$, and $\mathcal{E}(m,T)$ is the tail.
Expanding the kernel:
\begin{equation}\label{eq:kernelexp}
  \frac{(m+k_1)^s - m^s}{s}
  = m^{s-1}k_1\!\left(1 + O\!\left(\frac{|s|k_1}{m}\right)\right).
\end{equation}
For $m \geq k_1^2$ and $s$ near the saddle point ($|\Im(s)| \leq C/\log m$),
$|s| \approx \sigma_0 \approx 1$, so $|s|k_1/m \leq C/k_1 \to 0$.
The leading-term integral $\frac{k_1}{2\pi i}\int F_{k_1}(s)m^{s-1}ds$
yields $k_1\rho(u_1)(1+O(1/\log k_1))$ by the HT saddle-point analysis.
For $m \gg k_1\log n$ the correction $O(k_1/m)$ is $o(1)$, giving
$k_1\rho(u_1)(1+o(1)) \geq 1$ for large $n$.
The range $m \in [k_1^2, k_1\log n \cdot k_1] = [k_1^2, k_1^2\log n]$ is covered by the same argument (the correction is $O(k_1/k_1^2) = O(1/k_1) = o(1)$ for $m \geq k_1^2$).

This sketch glosses over three genuine technical obstacles, which
a complete proof would need to address.

\textbf{Obstacle~1} (Bounding the correction-term integral).
The correction $O(|s|k_1/m)$ in \eqref{eq:kernelexp} contributes
$O(k_1^2/m)$ times the integral $\int_{\sigma_0-iT}^{\sigma_0+iT}|s||F_{k_1}(s)|m^{\sigma_0-1}|ds|$.
This $s$-weighted integral is \emph{not} a direct consequence of
the Perron formula for $\Psi(m,k_1)$: inserting an extra factor $|s|^2/m$
into the Perron integrand requires bounding $|F_{k_1}(\sigma_0+it)|$
for $|t|$ up to $T$, and such vertical-strip bounds must be stated explicitly.

\textbf{Obstacle~2} (Tail uniformity in $m$).
The tail satisfies
\[
  \mathcal{E}(m,T) = O\!\left(\frac{(m+k_1)^{\sigma_0}F_{k_1}(\sigma_0)}{T}\right)
  \approx O\!\left(\frac{m\rho(u_1)}{T}\right).
\]
To achieve $\mathcal{E}(m,T) = o(k_1\rho(u_1))$ one needs $T \gg m/k_1$.
Since $\sigma_0 = 1 - (\log u_1)/\log k_1$ depends on $m$,
both the choice of $T$ and the resulting error must be controlled
\emph{uniformly} over all $m \in [k_1^2, n]$.
As $m/k_1$ ranges up to $n/k_1 \to \infty$, this requires
the saddle-point analysis to remain valid for $T$ growing without bound,
which is a non-trivial uniformity statement.

\textbf{Obstacle~3} (Contour validity for large $|\Im(s)|$).
The expansion \eqref{eq:kernelexp} and the bound on $F_{k_1}(\sigma_0+it)$
must hold uniformly for $|\Im(s)|$ up to $T = O(m/k_1) \to \infty$.
In the HT framework, $|F_y(\sigma_0+it)| \ll |F_y(\sigma_0)|e^{-c(\log y)|t|}$
for some $c > 0$ when $|t| \leq 1$; extending this to $|t| \to \infty$
as required here demands additional input on the distribution of $F_y(s)$
on tall rectangles, which is available (it follows from zero-free regions for
$F_y(s)$) but must be made explicit.

\begin{remark}[Summary: what an unconditional proof requires]\label{rem:obstacles}
Both routes are refinements of the HT saddle-point method, not new ideas.
Route~1 requires a second-order uniform expansion
$\Psi(x,y) = x\rho(u)(1 + c_1(u)/\log y + O((\log y)^{-2}))$
with $c_1(u)$ satisfying a Lipschitz bound, uniform in $x$.
Route~2 requires resolving Obstacles~1--3 above: bounding
$s$-weighted Perron integrals, controlling tails uniformly in $m$,
and extending saddle-point bounds to tall contours.
All necessary inputs---the Perron integral for $F_y(s)m^s$,
bounds on $|F_y(\sigma+it)|$, the Dickman function and its derivatives---are
available in principle via \cite{HT,Ten}, but the arguments have not been
assembled into a complete proof of Hypothesis~\ref{hyp:SI}.
\end{remark}

\section{Main Result}

\begin{theorem}\label{thm:main}
\begin{enumerate}
\item[\rm(a)] (Unconditional) For every fixed $\varepsilon > 0$ and all large $n$:
  $k(n) \geq \exp((\log n)^{1/2-\varepsilon})$.
\item[\rm(b)] (Conditional on Hypothesis~\ref{hyp:SI}) For every fixed $\varepsilon > 0$
  and all large $n$: $k(n) \leq \exp((\log n)^{1/2+\varepsilon})$.
\end{enumerate}
Hence, conditional on Hypothesis~\ref{hyp:SI}, $\log k(n) = (\log n)^{1/2+o(1)}$.
Hypothesis~\ref{hyp:SI} is a short-interval smooth number estimate
plausibly within reach of the Hildebrand--Tenenbaum method
but not verified directly from the published sources \cite{HT,Ten} here.
\end{theorem}

\begin{proof}
Part~(a) is Theorem~\ref{thm:lower}.
Part~(b) is Theorem~\ref{thm:upper}.
\end{proof}

\begin{thebibliography}{9}

\bibitem{HT}
  A.\ Hildebrand and G.\ Tenenbaum,
  \emph{Integers without large prime factors},
  J.\ Th\'eor.\ Nombres Bordeaux \textbf{5} (1993), 411--484.
\bibitem{Ten}
  G.\ Tenenbaum,
  \emph{Introduction to Analytic and Probabilistic Number Theory},
  Cambridge University Press, 3rd ed., 2015.

\end{thebibliography}

\end{document}
